\specialchapter{Historical Note}

(N.B. --- The Roman numerals refer to the bibliography at the end of this note.)

While the study of the connections between topology and measure theory goes back to the beginnings of the modern theory of functions of real variables, it is only very recently that integration in Hausdorff topological spaces has been thoroughly worked out in general fashion. Before setting out the history of the work that preceded the present synthesis, we review several stages in the evolution of the ideas regarding the relations between topology and measure.

For Lebesgue, it is only a question of integrating functions of one or several real variables. In 1913, Radon defined general measures on \(\mathbf{R}^n\) and the corresponding integrals; this theory is exposed in detail in the book (I) of Ch. de la Vallée Poussin and rests, in constant fashion, on the topological properties of Euclidean spaces. A little later, in 1915, Fréchet defined in (II, \(a\) ) `abstract' measures on a set equipped with a tribe, and the integrals with respect to these measures; he noted that one can establish in this way the principal results of the Lebesgue theory without the use of topological methods. He justified his undertaking with the following words, taken from the introduction of (II, \(b\) ), first part): « Que par exemple dans l'espace à une infinité de coordonnées où diverses applications de l'Analyse avaient conduit à diverses définitions non équivalentes d'une suite convergente, on remplace une de ces définitions par une autre, rien ne sera changé dans les propriétés des familles et fonctions additives d'ensembles dans ces espaces».(*) Fréchet's investigations were completed by Carathéodory, to whom is due an important theorem on the extension of a set function to a measure. The beginning of the book of Saks (III) gives a condensed exposition of this point of view.

The discovery of Haar measure on locally compact groups (cf.~the Historical Note for Chs. VII and VIII) and the numerous applications that it immediately received, then the works of Weil and Gelfand on Harmonic Analysis, led around 1940 to a profound modification of this point of view: in this kind of question, it is most convenient to regard a measure as a linear form on a space of continuous functions. This method obliges one to restrict oneself to compact or locally compact spaces, but this is not an inconvenience for nearly all of the applications: better yet, the introduction of Harmonic Analysis on p-adic groups and adèle groups by J. Tate and A. Weil has permitted a spectacular renewal of Analytic Number Theory.

It is from an entirely different direction that the need arises for enlarging this point of view by considering measures on non locally compact topological spaces: gradually, Probability Theory leads to the study of such spaces and furnishes numerous nontrivial examples. Perhaps the reason for the tardy influence of these developments on the theory of measure is to be found in the relative isolation of Probability Theory, which remained until recently on the fringes of the traditional mathematical disciplines.

